\chapter{Goal Setting}
\index{Goal Setting}

\section*{Definition and Overview}
Goal setting is the process of identifying specific, meaningful objectives and developing a plan to achieve them. In health care, goal setting is used to help people change behaviour, manage chronic conditions, recover from illness or injury, and improve wellbeing. Well-formed goals are specific, measurable, achievable, relevant, and time-bound (the SMART framework). Setting goals works because it turns vague intentions into concrete actions, provides motivation and direction, and allows progress to be tracked. Health professionals, including GPs, physiotherapists, psychologists, and exercise physiologists, routinely use goal setting to support behaviour change in areas such as physical activity, weight management, smoking cessation, and rehabilitation.

\section*{Epidemiology}
Goal setting is a core component of most effective behaviour-change interventions, and behaviour change is central to preventing and managing the leading causes of disease. In many countries, most adults have at least one modifiable health risk, such as insufficient physical activity, unhealthy diet, smoking, or excess alcohol, and more than half of adults live with a chronic condition. Structured programs that include goal setting have been shown to improve outcomes across weight loss, exercise adoption, diabetes self-management, cardiac rehabilitation, and mental health treatment, compared with advice alone.

\section*{Aetiology and Pathophysiology}
This chapter concerns behaviour change rather than a disease. The psychological mechanisms behind effective goal setting include:

\begin{itemize}
    \item \textbf{Clarity and focus:} specific goals direct attention and effort towards defined actions
    \item \textbf{Motivation:} meaningful goals generate commitment and persistence
    \item \textbf{Feedback:} measurable goals allow progress to be tracked, reinforcing effort
    \item \textbf{Self-efficacy:} achieving small goals builds confidence in the ability to change
    \item \textbf{Planning:} goals linked to an action plan ("when, where, how") are more likely to be achieved
    \item \textbf{Accountability:} sharing goals with a health professional or supporter increases follow-through
    \item \textbf{Adjustment:} realistic goals can be revised in response to progress and setbacks
\end{itemize}

\section*{Clinical Features}
Situations in which goal setting is particularly useful include:

\begin{itemize}
    \item Starting or increasing physical activity
    \item Losing weight or changing eating habits
    \item Quitting smoking, vaping, or reducing alcohol
    \item Managing a chronic condition such as diabetes, heart disease, or asthma
    \item Recovering from surgery, injury, or stroke
    \item Improving mental health, sleep, or stress management
    \item Reducing medication use or improving medicine adherence
    \item Returning to work or valued activities after illness
\end{itemize}

\section*{Differential Diagnosis}
Effective goal setting takes account of factors that might otherwise undermine it.

\begin{itemize}
    \item Readiness to change: goals work best when the person is motivated, not pressured
    \item Underlying physical or mental health problems that need treatment first
    \item Practical barriers, including cost, transport, time, and competing demands
    \item Social and environmental influences, including family and workplace culture
    \item Lack of knowledge or skills, which may need education before change
    \item Previous failed attempts, which can affect confidence
\end{itemize}

\section*{When to Seek Medical Care}
People with health conditions should set health goals in partnership with their care team. Before starting a new exercise program or making major dietary changes, people with chronic conditions, injuries, or concerns should discuss their plans with a health professional.

\textbf{Contact your local emergency services for:}
\begin{itemize}
    \item This chapter concerns behaviour change rather than emergencies; however, seek urgent care for any new chest pain, breathlessness, fainting, or other serious symptom during physical activity
    \item Mental health crises, including thoughts of self-harm, require immediate support
\end{itemize}

\section*{Diagnosis and Investigations}
Goal setting is individualised and requires understanding the person's starting point and context.

\begin{itemize}
    \item \textbf{Assessment:} current behaviour, health status, readiness to change, and confidence
    \item \textbf{Exploring values:} what matters to the person and why the goal is important
    \item \textbf{Identifying barriers and enablers:} practical, social, and psychological factors
    \item \textbf{Baseline measures:} weight, blood pressure, fitness, or other measures to track progress
    \item \textbf{Ongoing review:} regular check-ins to monitor progress and adjust goals
\end{itemize}

\section*{Management}
Goal setting is most effective when it is structured and supported.

\begin{itemize}
    \item \textbf{Use the SMART framework:} make goals Specific, Measurable, Achievable, Relevant, and Time-bound
    \item \textbf{Set one or two goals at a time:} focusing on a small number of goals improves success
    \item \textbf{Make an action plan:} decide when, where, and how the goal will be actioned
    \item \textbf{Start small:} early success builds confidence and momentum
    \item \textbf{Write goals down:} written and shared goals are more likely to be achieved
    \item \textbf{Track progress:} diaries, apps, or checklists make progress visible
    \item \textbf{Anticipate setbacks:} plan how to respond to slips so they do not derail the goal
    \item \textbf{Review and adjust:} revise goals in response to progress, life changes, and feedback
    \item \textbf{Celebrate success:} acknowledge achievements and set new goals as old ones are met
\end{itemize}

\section*{Complications}
\begin{itemize}
    \item Unrealistic or overly ambitious goals that lead to frustration and giving up
    \item Vague goals that cannot be measured or acted on
    \item Too many goals at once, which dilutes focus
    \item All-or-nothing thinking, where a single slip is treated as failure
    \item Goals set by others rather than by the person, which reduces commitment
    \item Failure to adapt goals when circumstances change
    \item Loss of motivation if progress is not reviewed or celebrated
\end{itemize}

\section*{Prognosis and Outcomes}
People who set specific, realistic, and self-chosen goals are substantially more likely to change behaviour than those who receive advice alone. Goal-based programs produce meaningful improvements in physical activity, weight, blood pressure, blood glucose, and psychological wellbeing, and these gains are more likely to persist when goals are reviewed over time. The skills of goal setting are transferable: people who learn to set and achieve goals in one area of health often apply the approach to others, building a foundation of self-management that supports lifelong health.

\section*{Prevention and Self-Care}
\begin{itemize}
    \item Choose goals that genuinely matter to you, not just to others
    \item Make each goal specific and measurable, such as "walk 30 minutes on five days"
    \item Set a timeframe for achieving each goal
    \item Plan exactly when and where you will act on the goal
    \item Start with one small change and build on success
    \item Keep a simple record of progress
    \item Ask a friend, family member, or health professional to support and check in with you
    \item Be kind to yourself after setbacks and simply resume the plan
    \item Review your goals regularly and adjust them as your life changes
\end{itemize}

\section*{Sources and Further Reading}
The following reputable sources provide further information for healthcare students and patients seeking reliable, current advice about goal setting for health.

\begin{itemize}
    \item Healthdirect Australia. \textit{Setting health goals that work.} https://www.healthdirect.gov.au/
    \item Heart Foundation. \textit{Healthy living goals and action plans.} https://www.heartfoundation.org.au/
    \item Exercise and Sports Science Australia. \textit{Setting exercise goals with a professional.} https://www.essa.org.au/
    \item Australian Government Department of Health and Aged Care. \textit{Healthy living and behaviour change.} https://www.health.gov.au/
    \item NHS. \textit{Setting goals for a healthier lifestyle.} https://www.nhs.uk/
    \item World Health Organization. \textit{Behaviour change and health promotion.} https://www.who.int/
\end{itemize}
